\documentclass[12pt,a4paper]{article}
\usepackage[margin=1in]{geometry}
\usepackage[T1]{fontenc}
\usepackage{setspace}
\usepackage{amsmath,amssymb}
\usepackage{booktabs}
\usepackage{longtable}
\usepackage{xcolor}
\usepackage{hyperref}
\usepackage{fancyhdr}

\setstretch{1.18}
\setlength{\parskip}{0.45em}
\setlength{\parindent}{0pt}
\setlength{\headheight}{15pt}
\hypersetup{colorlinks=true,linkcolor=blue,citecolor=blue,urlcolor=blue}

\pagestyle{fancy}
\fancyhf{}
\lhead{CST428 - Blockchain Technologies}
\rhead{Module 4 Master Notes}
\cfoot{\thepage}

\newcommand{\examtip}[1]{%
\vspace{0.35em}
\noindent\fcolorbox{black}{gray!12}{\parbox{0.965\linewidth}{\textbf{Exam Tip:} #1}}%
\vspace{0.35em}
}

\begin{document}

\section{Module 4 Topic Map From Previous Questions}

From May 2024, June 2023, and October 2023 papers, Module 4 repeatedly asks:

\begin{itemize}
\item Smart contracts: definition, properties, and practical meaning.
\item Oracles: types, oracle problem, and smart contract data flow.
\item DApps: definition, architecture, and design flow with diagram.
\item Sector applications: government, finance, healthcare, supply chain.
\item Emerging combinations: blockchain with AI and cloud computing.
\end{itemize}

\section{Core Theory (Explained Clearly)}

\subsection{Smart Contracts: The Core Idea}
A smart contract is a program stored on a blockchain that runs automatically when predefined conditions are satisfied.

Simple way to think:
\begin{itemize}
\item Traditional contract: written on paper, humans enforce it.
\item Smart contract: written in code, network enforces it.
\end{itemize}

Key properties:
\begin{itemize}
\item \textbf{Automatic execution:} no manual follow-up needed after trigger condition.
\item \textbf{Immutable logic:} after deployment, rules are hard to change.
\item \textbf{Transparent behavior:} anyone can verify what contract will do.
\item \textbf{Deterministic output:} same input gives same output on all nodes.
\end{itemize}

\examtip{In 3-mark questions, one strong definition plus 4 properties is enough for full marks.}

\subsection{What Is an Oracle and Why Is It Needed?}
Blockchains are excellent at verifying on-chain data, but they cannot directly read external web APIs.
Oracles solve this by bringing trusted external data into smart contracts.

Examples of oracle data:
\begin{itemize}
\item weather data for crop insurance payout,
\item market prices for DeFi lending,
\item shipment status for supply chain release,
\item election result feeds for governance triggers.
\end{itemize}

\subsubsection*{Common Oracle Types}
\begin{itemize}
\item \textbf{Software oracle:} fetches digital data from web services.
\item \textbf{Hardware oracle:} reads sensor or IoT device data.
\item \textbf{Inbound oracle:} external data to blockchain.
\item \textbf{Outbound oracle:} blockchain event to external system.
\item \textbf{Centralized oracle:} single provider, simple but risky.
\item \textbf{Decentralized oracle:} multiple providers, safer through aggregation.
\end{itemize}

Useful formula used by decentralized oracle systems:
\[
\text{Final Price} = \mathrm{median}(p_1, p_2, \dots, p_n)
\]
Median helps reduce manipulation by outlier values.

\begin{center}
\framebox[0.95\linewidth][c]{
\begin{minipage}{0.90\linewidth}
\vspace{0.1cm}
\centerline{\textbf{Generic Oracle Data Flow}}
\centerline{Smart Contract requests data}
\centerline{$\Downarrow$}
\centerline{Oracle network fetches from APIs/sensors}
\centerline{$\Downarrow$}
\centerline{Responses are validated and aggregated}
\centerline{$\Downarrow$}
\centerline{Trusted value posted back on-chain}
\centerline{$\Downarrow$}
\centerline{Smart contract executes rule}
\vspace{0.1cm}
\end{minipage}}
\end{center}

\subsection{DApps: What Makes Them Different}
A DApp (Decentralized Application) is an application where backend logic runs on smart contracts instead of a centralized server.

\subsubsection*{DApp Building Blocks}
\begin{itemize}
\item \textbf{Frontend:} web or mobile interface.
\item \textbf{Wallet layer:} identity and signing (for example, MetaMask).
\item \textbf{Web3 library:} bridge between UI and blockchain.
\item \textbf{Smart contract backend:} business rules on chain.
\item \textbf{Node/RPC service:} path to read/write blockchain state.
\item \textbf{Optional off-chain storage:} large files in IPFS/Arweave.
\end{itemize}

\begin{center}
\framebox[0.95\linewidth][c]{
\begin{minipage}{0.90\linewidth}
\vspace{0.1cm}
\centerline{\textbf{Basic DApp Architecture}}
\centerline{User Interface (Web/Mobile)}
\centerline{$\Downarrow$}
\centerline{Wallet + Web3 Library}
\centerline{$\Downarrow$}
\centerline{RPC Node}
\centerline{$\Downarrow$}
\centerline{Smart Contract (On-chain Logic)}
\centerline{$\Downarrow$}
\centerline{Blockchain State + Events}
\vspace{0.1cm}
\end{minipage}}
\end{center}

Simple throughput idea for DApps:
\[
\text{User Throughput} \approx \frac{\text{Network TPS}}{\text{Average Calls Per User Action}}
\]
If one user action needs multiple contract calls, perceived speed decreases.

\subsection{Where Module 4 Is Used in Real Life}

\subsubsection*{Government Services}
\begin{itemize}
\item tamper-evident land records,
\item transparent subsidy disbursal,
\item auditable procurement logs,
\item verifiable digital certificates.
\end{itemize}

\subsubsection*{Finance Use Cases}
\begin{itemize}
\item instant settlement without many intermediaries,
\item programmable escrow,
\item tokenized assets and fractional ownership,
\item transparent audit trails for compliance.
\end{itemize}

\subsubsection*{Healthcare Use Cases}
\begin{itemize}
\item patient-consent-controlled data sharing,
\item immutable access logs,
\item anti-counterfeit medicine tracking,
\item inter-hospital record interoperability.
\end{itemize}

\subsubsection*{Supply Chain Use Cases}
\begin{itemize}
\item source-to-shelf product traceability,
\item cold-chain sensor-backed verification,
\item delayed-payment automation via smart contracts,
\item dispute reduction through shared immutable records.
\end{itemize}

\subsubsection*{Blockchain + AI and Cloud}
\begin{itemize}
\item AI: blockchain gives verifiable data lineage for model trust.
\item AI: smart contracts can automate model access and usage payment.
\item Cloud: blockchain audit trails improve integrity and compliance checks.
\item Cloud: multi-organization workflows can share trust without one central owner.
\end{itemize}

\section{Solved Question Papers}

\subsection{Part A: 3-Mark Questions}

\subsubsection*{Q1. What are smart contracts? Mention their properties.}
\begin{itemize}
\item A smart contract is a blockchain-stored program that auto-executes agreed rules.
\item It removes manual enforcement after deployment.
\item It is immutable in normal operation, so logic tampering is difficult.
\item It is transparent because rule behavior can be audited.
\item It is deterministic: same input gives same output on all nodes.
\item It supports trust-minimized transactions between unknown parties.
\end{itemize}

\subsubsection*{Q2. Explain different types of oracles.}
\begin{itemize}
\item Software oracle collects digital data from APIs and websites.
\item Hardware oracle collects physical-world data from sensors and devices.
\item Inbound oracle sends off-chain data to smart contracts.
\item Outbound oracle sends blockchain events to external systems.
\item Centralized oracle uses one source and has single-point risk.
\item Decentralized oracle aggregates multiple sources to improve trust.
\end{itemize}

\subsubsection*{Q3. What are DApps?}
\begin{itemize}
\item DApps are applications whose backend logic runs on smart contracts.
\item Users interact using wallets instead of username-password models.
\item Transaction authorization is done by digital signatures.
\item Business logic becomes transparent and verifiable.
\item Censorship resistance is stronger than many centralized apps.
\item They still use frontend UI, but trust anchor is on-chain logic.
\end{itemize}

\subsection{Part B: 14-Mark Questions}

\subsubsection*{Q1. Discuss oracles in blockchain ecosystem. Explain generic data flow from smart contract to oracle.}
\begin{itemize}
\item Smart contracts cannot directly call external APIs due to determinism and consensus limits.
\item Oracles bridge this gap by fetching and delivering external data.
\item Typical use areas are DeFi price feeds, insurance triggers, and supply-chain status updates.
\item The smart contract raises a data request event.
\item Oracle nodes monitor request events and fetch requested off-chain data.
\item Data is validated, normalized, and often aggregated.
\item Aggregation can use median to reduce manipulation impact.
\item Oracle network publishes verified result back on-chain.
\item Smart contract reads returned value and executes programmed action.
\item Decentralized oracle design reduces single-point failure and censorship risk.
\item Oracle integrity is still a system risk area called the oracle problem.
\item Reputation, staking, and multi-source verification are common protections.
\item Secure oracle integration is mandatory for production-grade smart contracts.
\item[]
\vspace{0.25em}
\noindent\framebox[0.97\linewidth][c]{
\begin{minipage}{0.93\linewidth}
\vspace{0.08cm}
\centerline{\textbf{Oracle Flow Diagram}}
\centerline{Contract Request $\rightarrow$ Oracle Fetch $\rightarrow$ Validate/Aggregate}
\centerline{$\Downarrow$}
\centerline{On-chain Response $\rightarrow$ Contract Action}
\vspace{0.08cm}
\end{minipage}}
\end{itemize}

\subsubsection*{Q2. Illustrate ways blockchain can be employed in government services.}
\begin{itemize}
\item Land and property records can be made tamper-evident and publicly verifiable.
\item Certificate issuance can use verifiable credentials with instant authenticity checks.
\item Subsidy distribution can use transparent eligibility and payment audit trails.
\item Public procurement can be logged with immutable milestones to reduce corruption.
\item Identity systems can move toward citizen-controlled verifiable identity claims.
\item Voting-related modules can improve transparency for selected workflows.
\item Cross-border document attestation can become faster and less fraud-prone.
\item Smart contracts can automate conditional grants and compliance checks.
\item Permissioned blockchain can preserve privacy while enabling traceability.
\item Adoption success depends on legal process redesign and governance policy.
\end{itemize}

\subsubsection*{Q3. Illustrate with a use case: blockchain application in finance sector.}
\begin{itemize}
\item Consider cross-border remittance between two banks and a settlement partner.
\item Traditional process involves multiple intermediaries and delayed reconciliation.
\item With blockchain, transaction status is shared on a common ledger.
\item Smart contract can verify KYC status and transaction rules automatically.
\item Payment release can be conditional on compliance checks and confirmations.
\item Settlement can occur faster with less manual matching.
\item Audit trail is immutable and easier for regulators to inspect.
\item Dispute resolution time decreases due to common data visibility.
\item Programmable escrow reduces counterparty risk in trade finance.
\item Operational cost can reduce by minimizing duplicate reconciliation work.
\item Risk controls must include oracle quality, compliance, and private key safety.
\end{itemize}

\subsubsection*{Q4. What are DApps? Outline/basic architecture and design of a DApp with neat diagram.}
\begin{itemize}
\item A DApp is an app where core logic is executed by smart contracts.
\item The UI can look similar to normal apps but trust model is different.
\item User identity and authorization are wallet-based via signatures.
\item Web3/Ethers layer translates UI actions to blockchain calls.
\item Read operations query current chain state through RPC endpoints.
\item Write operations create signed transactions and wait for confirmations.
\item Contract events are used to update UI state predictably.
\item Off-chain storage is used for large media to reduce blockchain cost.
\item Security review includes contract audit and frontend signing safety.
\item Gas optimization is a design requirement for user affordability.
\item Error handling must include transaction pending, revert, and timeout states.
\item Production DApps need monitoring for RPC health and oracle consistency.
\item[]
\vspace{0.25em}
\noindent\framebox[0.97\linewidth][c]{
\begin{minipage}{0.93\linewidth}
\vspace{0.08cm}
\centerline{\textbf{DApp Design Diagram}}
\centerline{UI $\leftrightarrow$ Wallet $\leftrightarrow$ RPC Node}
\centerline{$\Downarrow$}
\centerline{Smart Contract $\leftrightarrow$ On-chain State}
\centerline{$\Downarrow$}
\centerline{Optional Oracle + IPFS Layers}
\vspace{0.08cm}
\end{minipage}}
\end{itemize}

\subsubsection*{Q5. Explain impact of blockchain technology on artificial intelligence.}
\begin{itemize}
\item AI quality depends heavily on data quality and data trust.
\item Blockchain can provide verifiable data lineage and timestamped provenance.
\item This helps teams prove where training data came from.
\item Model sharing can be controlled by smart contracts and access policies.
\item Payments for model usage can be automated per inference request.
\item Federated AI collaborations can share results without full data exposure.
\item Tamper-evident logs help with model governance and auditability.
\item On-chain storage is expensive, so hybrid design is necessary.
\item Privacy, latency, and regulation remain practical constraints.
\item Best architecture combines blockchain trust layer with off-chain AI compute.
\item Use blockchain where auditability and multi-party trust are critical.
\end{itemize}

\subsubsection*{Q6. Explain how blockchain can impact healthcare sector.}
\begin{itemize}
\item Patient records can be shared with explicit and revocable consent.
\item Access logs become immutable, improving accountability.
\item Drug supply chain can be tracked to detect counterfeit medicines.
\item Insurance claims can use smart contracts for faster adjudication.
\item Clinical trial records can be timestamped for integrity.
\item Interoperability improves when multiple hospitals trust same ledger.
\item Patients can maintain stronger control over who accesses what data.
\item Data integrity helps reduce accidental record tampering.
\item Encryption and permissioning are mandatory for privacy compliance.
\item Off-chain storage with on-chain proofs is practical for large files.
\item Governance and legal compliance determine real-world deployment success.
\end{itemize}

\subsubsection*{Q7. Illustrate how blockchain can be used in supply chain management.}
\begin{itemize}
\item Every handoff event can be logged as a verifiable transaction.
\item Product origin can be checked quickly using immutable history.
\item Sensor-backed data (temperature, location) can trigger contract conditions.
\item Payments can be released automatically when delivery milestones are met.
\item Shared ledger reduces disputes between suppliers and buyers.
\item Counterfeit insertion becomes harder due to chain-of-custody proofs.
\item Recall operations become faster by tracing affected batches.
\item Multi-party planning can improve when status data is synchronized.
\item Smart contracts reduce manual paperwork in shipment workflows.
\item Permissioned blockchain helps preserve commercial confidentiality.
\item Integration with ERP systems is required for enterprise adoption.
\item Overall impact is improved trust, traceability, and operational speed.
\end{itemize}

\subsubsection*{Q8. Explain benefits and applications of integrating blockchain in cloud computing.}
\begin{itemize}
\item Cloud environments need strong auditability across many actors.
\item Blockchain can provide immutable logs of critical cloud events.
\item Multi-cloud workflows can coordinate trust without one dominant operator.
\item Smart contracts can automate SLA checks and penalty enforcement.
\item Identity and access events can be made tamper-evident.
\item Billing disputes can reduce through transparent metering evidence.
\item Data integrity checks can use hash anchoring on chain.
\item Federated organizations can collaborate with auditable governance rules.
\item Throughput and latency constraints require selective on-chain use.
\item Sensitive data should remain encrypted and mostly off-chain.
\item Blockchain should complement cloud, not replace core cloud services.
\end{itemize}

\end{document}
